%! AP Statistics — Unit 4, Lesson 4.2 — Student Guided Notes
\documentclass[10pt]{article}
\usepackage{apstats-article}
\usepackage{apstats-boxes}

\begin{document}

\pageheader{Unit 4, Lesson 4.2}{Guided Notes}
\namedateperiod

\vspace{0.12in}

\begin{objectivebox}
\small By the end of this lesson, I will be able to\dots\\[4pt]
\writeline\\[1pt]
\writeline\\[1pt]
\writeline
\end{objectivebox}

\vspace{0.1in}

\begin{vocabbox}
\termblanklong{Simulation}
\termblanklong{Trial}
\termblanklong{Outcome}
\termblanklong{Event}
\termblanklong{Relative frequency (empirical probability)}
\termblanklong{Law of Large Numbers}
\end{vocabbox}

\vspace{0.1in}

\begin{hookbox}
\small
A thumbtack is dropped 50 times. Tally the results and compute running relative frequency.

\vspace{6pt}
\renewcommand{\arraystretch}{1.6}
\begin{tabularx}{\linewidth}{@{} l c c c c c @{}}
After $n$ drops: & 5 & 10 & 20 & 35 & 50 \\
\hline
Point-up count & \blank{1cm} & \blank{1cm} & \blank{1cm} & \blank{1cm} & \blank{1cm} \\
Relative freq. & \blank{1cm} & \blank{1cm} & \blank{1cm} & \blank{1cm} & \blank{1cm} \\
\end{tabularx}

\vspace{6pt}
What do you notice about the relative frequency as $n$ increases? \writeline
\end{hookbox}

\vspace{0.1in}

\begin{notesbox}{Simulation Design Framework}
\small

\textbf{Key Idea (UNC-2.A.4):} Simulation is a way to model random events such that
simulated outcomes \blank{5cm} real-world outcomes. All possible outcomes are
associated with a value determined by \blank{2cm}.

\vspace{10pt}
\textbf{Steps to Design a Simulation:}
\begin{enumerate}[leftmargin=1.5em, itemsep=6pt]
  \item \textbf{Identify} the random process: \blank{8cm}
  \item \textbf{Assign} values to outcomes: \blank{8cm}
  \item \textbf{Choose} a chance device: \blank{8cm}
  \item \textbf{Run} trials and record outcomes: \blank{8cm}
  \item \textbf{Compute} relative frequency: $\hat{p} = \dfrac{\blank{3cm}}{\blank{3cm}}$
\end{enumerate}

\vspace{10pt}
\textbf{Example:} A genetic cross. Each parent is a carrier (one dominant, one recessive allele).
Each offspring independently receives one allele from each parent.
\begin{itemize}[itemsep=4pt]
  \item Chance device: \blank{5cm}
  \item Assign: \blank{7cm}
  \item One trial = \blank{5cm}
  \item Event of interest: offspring has two recessive alleles (affected)
\end{itemize}
\end{notesbox}

\vspace{0.1in}

\begin{notesbox}{The Law of Large Numbers}
\small
\textbf{UNC-2.A.5:} The relative frequency of an outcome in simulated data can be used
to \blank{5cm} the probability of that outcome.

\vspace{8pt}
\textbf{UNC-2.A.6 --- Law of Large Numbers:}\\
As the number of trials \blank{3cm}, the simulated (empirical) probability tends to
get \blank{4cm} to the \blank{4cm} probability.

\vspace{8pt}
\textbf{What this means in practice:}
\begin{itemize}[itemsep=4pt]
  \item With \blank{2cm} trials, relative frequency can vary a lot.
  \item With \blank{2cm} trials, relative frequency stabilizes near the true value.
  \item A few unlucky trials early on \blank{4cm} (do/does not) mean the estimate is wrong;
        it just takes more trials to stabilize.
\end{itemize}
\end{notesbox}

\vspace{0.1in}

\begin{practicebox}
\small
A free-throw shooter makes 72\% of attempts. Use the random number table
below to simulate 20 shots. Let digits 00--71 represent a make and 72--99
represent a miss.

\vspace{6pt}
{\small\texttt{39 48 54 70 29 73 41 61 85 33 \quad 17 56 88 12 65 09 42 99 77 23}}

\vspace{8pt}
\begin{enumerate}[leftmargin=1.5em, itemsep=8pt]
  \item List the outcome (make/miss) for each of the 20 two-digit values above.
        \vspace{20pt}
  \item How many of the 20 shots were makes? \blank{2cm}
  \item What is the empirical probability (relative frequency) of a make?
        $\hat{p} = \blank{3cm}$
  \item How does this compare to the true probability of 0.72?
        \writeline
\end{enumerate}
\end{practicebox}

\end{document}
